\documentclass[11pt]{amsart}
\usepackage[T1]{fontenc}
\usepackage{lmodern,microtype,amsmath,amssymb,amsthm,mathtools,enumitem,booktabs}
\usepackage[margin=1.1in]{geometry}
\usepackage[hidelinks]{hyperref}
\usepackage{xurl}

\newtheorem{theorem}{Theorem}[section]
\newtheorem{lemma}[theorem]{Lemma}
\newtheorem{proposition}[theorem]{Proposition}
\newtheorem{corollary}[theorem]{Corollary}
\theoremstyle{remark}
\newtheorem{remark}[theorem]{Remark}
\numberwithin{equation}{section}

\newcommand{\EE}{\mathbb{E}}
\newcommand{\RR}{\mathbb{R}}
\DeclareMathOperator{\Hess}{Hess}

\hypersetup{pdftitle={Gradient flow for factored scalar linear attention},pdfauthor={Yitzchak Shmalo}}

\title[Gradient flow for factored scalar linear attention]{Gradient flow for deep factored scalar linear attention learns multi-step gradient descent}
\author{Yitzchak Shmalo}
\address{Einstein Institute of Mathematics, The Hebrew University of Jerusalem, Jerusalem, Israel}
\email{yitzchak.shmalo@gmail.com}
\date{November 2026}
\subjclass[2020]{68T07, 37N40, 33C45, 62J05}
\keywords{in-context learning, linear attention, gradient flow, strict saddle, Laguerre polynomials}

\begin{document}
\begin{abstract}
We study in-context learning of scalar linear regression by $L$ untied linear-attention layers, each with a trainable key scalar $k_i$ and query scalar $q_i$. For every depth $L$ and prompt length $n$, population gradient flow on the $2L$ factors is bounded and converges to a critical point, every critical point that is not a global minimizer has a negative Hessian eigenvalue, and the flow reaches a global minimizer from all initializations outside a Lebesgue null set, for instance almost surely under Gaussian initialization. The trained network performs $L$ gradient descent steps on the prompt's least-squares loss, with step sizes $k_iq_i$ the reciprocals of the zeros of a rescaled Laguerre polynomial $L_L^{(n/2)}$, and its risk is $2^LL!/\prod_{j=1}^{L}(n+2j)$. Initializations with $k_i=-q_i$ in some layer converge to a critical point that is not a minimizer, so the null set is not empty.
\end{abstract}
\maketitle

\section{Introduction}\label{sec:intro}
Trained transformers can solve linear regression problems presented in their prompt, and on such tasks their predictions are close to those of least squares~\cite{Garg}. Von Oswald et al.~\cite{vonOswald} showed that linear self-attention can implement gradient descent on the in-context loss, and Aky\"urek et al.~\cite{Akyurek} constructed transformer implementations of gradient-based and closed-form regression algorithms; both works also compared trained models with these constructions. Ahn, Cheng, Daneshmand and Sra~\cite{Ahn} described the loss landscape of linear transformers on random linear regression: for one layer the global minimum is one step of preconditioned gradient descent, and for a network of depth $L$, that is with $L$ attention layers, certain critical points correspond to $L$ such steps. Zhang, Frei and Bartlett~\cite{ZFB} proved that gradient flow on a single linear self-attention layer, from a suitable random initialization, converges to a global minimizer. Bordelon, Letey and Pehlevan~\cite{BLP} studied the dynamics of untied linear-attention layers in a high-dimensional limit, from equal isotropic or balanced small initializations; Theorem~\ref{thm:main} below holds at every finite prompt length, for scalar data, from almost every initialization.

For looped layers, which share one set of weights, Gatmiry, Saunshi, Reddi, Jegelka and Kumar~\cite{Gatmiry} prove that gradient flow converges and that its limit performs multi-step preconditioned gradient descent; their proof rests on a gradient dominance condition for the looped risk. We consider $L$ untied layers, that is, layers with separate weights, and scalar data. Layer $i$ has a key scalar $k_i$ and a query scalar $q_i$, called the factors, and the value map and the mask are fixed as in~\cite{Gatmiry}. Write $\theta=(k_1,q_1,\dots,k_L,q_L)$, let $n$ be the prompt length, and let $S$ be the empirical second moment of the prompt. The network output is the result of $L$ gradient descent steps on the prompt loss with step sizes $a_i=k_iq_i$, and the risk is
\[
f(\theta)=F(a)=\EE\Bigl[\,\prod_{i=1}^{L}(1-a_iS)^2\Bigr].
\]
No gradient dominance inequality holds for $f$ on $\RR^{2L}$: the origin is a critical point with $f(0)=1$, above the minimum. Moreover $f$ is constant along $(k_i,q_i)\mapsto(\lambda k_i,q_i/\lambda)$, so its sublevel sets are unbounded and its minimizers are not isolated, and a trajectory that starts with $k_i=-q_i$ in some layer converges to a critical point that is not a minimizer (Remark~\ref{rem:exceptional}). Theorem~\ref{thm:main} gives the counterpart of the result of~\cite{Gatmiry} for untied layers and scalar data: the flow is bounded and converges to a critical point, every critical point that is not a global minimizer is a strict saddle, and from almost every initialization the limit is a global minimizer. For deep linear networks, Bah, Rauhut, Terstiege and Westdickenberg~\cite{BRTW} showed that gradient flow on the factors is a Riemannian gradient flow of the product and converges to a critical point. From almost every initialization their flow converges to a global minimum on a manifold of matrices of some fixed rank. Here $F$ is not convex in $a$, and the reduction below goes through the coefficients of a polynomial instead.

The critical points are found in the coefficients of the polynomial $\prod_i(1-a_is)$. In these coefficients $F$ is a strictly convex quadratic, and the map from $a$ to the coefficients has invertible Jacobian wherever the $a_i$ are distinct. Hence a critical point of $F$ with distinct coordinates is a global minimizer. At a point with $a_i=a_j=\gamma$, write $\delta_i$ for the $i$-th coordinate vector and $R(s)=\prod_{l\ne i,j}(1-a_ls)$. Then
\[
F\bigl(a+t(\delta_i-\delta_j)\bigr)-F(a)=-2t^2\,\EE\bigl[S^2(1-\gamma S)^2R(S)^2\bigr]+t^4\,\EE\bigl[S^4R(S)^2\bigr],
\]
so the Hessian of $F$ is negative in the direction $\delta_i-\delta_j$ (Lemma~\ref{lem:Fland}). A critical point of $f$ that is not a global minimizer has two equal step sizes or a layer with $k_i=q_i=0$, and in both cases the Hessian of $f$ has a negative eigenvalue (Lemma~\ref{lem:saddle}).

The flow is bounded because each $k_i^2-q_i^2$ is conserved, as the balancedness of layers is for deep linear and homogeneous networks~\cite{ACH,DHL}:
\[
(k_i^2+q_i^2)^2=(k_i^2-q_i^2)^2+4a_i^2,
\]
and $a$ stays in a compact sublevel set of $F$ (Lemma~\ref{lem:bounded}). A bounded trajectory of the gradient flow of a polynomial converges to a single point by the theorem of \L{}ojasiewicz: an integral curve of the gradient flow of an analytic function that has an accumulation point has a unique limit, a consequence of his gradient inequality~\cite{ColdingMinicozzi}. The theorem of de~Oliveira, Siami and Sontag~\cite{OSS}, which allows a continuum of critical points, then shows that the initial points whose trajectories converge to a critical point other than a global minimizer form a Lebesgue null set.

The minimizer of $F$ is explicit. The optimal polynomial $\prod_i(1-a_is)$ is $L_L^{(n/2)}(ns/2)/L_L^{(n/2)}(0)$, where $L_L^{(n/2)}$ is the Laguerre polynomial~\cite{Szego}, as in the average-case analysis of Pedregosa and Scieur~\cite{PS}, who give the exponential case $n=2$, and of Cunha, Gidel, Pedregosa, Scieur and Paquette~\cite{CGPSP}, who treat Gamma measures; the optimal step sizes are the reciprocals of its zeros, and the minimal risk is $2^LL!/\prod_{j=1}^{L}(n+2j)$. The data are one-dimensional and noiseless, only the key and query scalars are trained, attention is linear, and training is by population gradient flow. The exponent in the \L{}ojasiewicz inequality is not explicit, and the rate of convergence is open.

\section{Setting and main result}\label{sec:setting}
Fix integers $n,L\ge1$. Let $x_1,\dots,x_n,x_q,w$ be independent standard normal variables and $y_j=wx_j$. Put
\[
Z_0=\begin{pmatrix}x_1&\cdots&x_n&x_q\\ y_1&\cdots&y_n&0\end{pmatrix},\qquad
M=\operatorname{diag}(1,\dots,1,0),\qquad P=\operatorname{diag}(0,1),
\]
and for $\theta=(k_1,q_1,\dots,k_L,q_L)\in\RR^{2L}$ let $K_i=\operatorname{diag}(k_i,0)$ and $Q_i=\operatorname{diag}(q_i,0)$. The layers are
\begin{equation}\label{eq:layer}
Z_i=Z_{i-1}-\frac1n\,PZ_{i-1}M\,(K_iZ_{i-1})^{\top}(Q_iZ_{i-1}),\qquad i=1,\dots,L,
\end{equation}
and the prediction for $y_q=wx_q$ is minus the bottom right entry of $Z_L$. The risk $f(\theta)$ is the expected squared prediction error. Write
\[
S=\frac1n\sum_{j=1}^{n}x_j^2,\qquad \alpha=\frac n2,\qquad (z)_j=z(z+1)\cdots(z+j-1),\quad (z)_0=1.
\]

Here $n$ is the prompt length, $L$ the depth, $Z_0$ the prompt with the query $x_q$ in its last column, $M$ the mask and $P$ the value map. A critical point of $f$ is a \emph{strict saddle} if the Hessian of $f$ there has a negative eigenvalue.

\begin{theorem}\label{thm:main}
Let $\theta(t)$ solve $\theta'=-\nabla f(\theta)$ with $\theta(0)\in\RR^{2L}$.
\begin{enumerate}[label=\textup{(\roman*)}]
\item The solution exists for all $t\ge0$, is bounded, and converges to a critical point of $f$.
\item Every critical point of $f$ that is not a global minimizer is a strict saddle.
\item The set $\mathcal N$ of initial points $\theta(0)$ whose solution does not converge to a global minimizer is a Lebesgue null set. In particular the solution converges to a global minimizer almost surely when $\theta(0)$ has a density, for example when the $2L$ factors are independent Gaussians.
\item At a global minimizer the products $a_i=k_iq_i$ are, in some order, the reciprocals of the $L$ zeros of
\begin{equation}\label{eq:rstar}
r_*(s)=\sum_{j=0}^{L}(-1)^j\binom Lj\frac{\alpha^js^j}{(\alpha+1)_j}=\frac{L!}{(\alpha+1)_L}\,L_L^{(\alpha)}(\alpha s),
\end{equation}
which are positive and simple, and the network predicts with $L$ gradient descent steps from zero on the prompt loss $\ell(v)=(2n)^{-1}\sum_j(vx_j-y_j)^2$ with these step sizes. The minimal risk is
\begin{equation}\label{eq:risk}
f_*=\frac{L!}{(\alpha+1)_L}=\frac{2^LL!}{\prod_{j=1}^{L}(n+2j)}.
\end{equation}
\end{enumerate}
\end{theorem}

Here $L_L^{(\alpha)}$ is the Laguerre polynomial. For $L=1$ the optimal step is $n/(n+2)$ and the risk is $2/(n+2)$; for $L=2$ the steps are $\frac n{n+2}\bigl(1\pm\sqrt{2/(n+4)}\bigr)$ and the risk is $8/((n+2)(n+4))$. The Bayes risk is zero, since $w=y_1/x_1$ almost surely, so~\eqref{eq:risk} is also the excess risk; at fixed $L$ it is asymptotic to $2^LL!\,n^{-L}$.

\begin{remark}\label{rem:exceptional}
The set $\mathcal N$ in~(iii) is not empty. Write $f(\theta)=F(a)$ with $a_i=k_iq_i$, as in~\eqref{eq:F} below. If $k_i(0)=-q_i(0)$ for some $i$, then $(k_i+q_i)'=-(k_i+q_i)\partial_{a_i}F$, so $k_i+q_i=0$ for all time and $a_i=-k_i^2\le0$, while every global minimizer has all $a_i>0$. For $L=1$ the flow on this line is $k'=-2k(1+m_2k^2)$ with $m_2=(\alpha+1)/\alpha$, so the factors go to zero and the risk to $1>2/(n+2)$. A law supported on such a hyperplane can be atomless, so the hypothesis that $\theta(0)$ has a density cannot be weakened to atomlessness.
\end{remark}

\section{Reduction to the products}\label{sec:reduction}
Since $K_i^{\top}Q_i=\operatorname{diag}(a_i,0)$, the layer~\eqref{eq:layer} depends on $\theta$ only through $a_i=k_iq_i$. The top row of $Z_i$ never changes because the first row of $P$ is zero. Suppose the bottom row of $Z_{i-1}$ is $((w-v)x_1,\dots,(w-v)x_n,-vx_q)$. The masked attention subtracts $a_iS(w-v)x_j$ from column $j$, the query column included, so the bottom row of $Z_i$ has the same form with $v$ replaced by $v+a_iS(w-v)=v-a_i\ell'(v)$. Starting from $v=0$, the prediction is the $L$-th gradient descent iterate on $\ell$ with steps $a_1,\dots,a_L$, and the prediction error is $wx_q\prod_i(1-a_iS)$. Integrating out $w$ and $x_q$ gives
\begin{equation}\label{eq:F}
f(\theta)=F(a),\qquad F(a)=\EE\bigl[r_a(S)^2\bigr],\qquad r_a(s)=\prod_{i=1}^{L}(1-a_is).
\end{equation}
The variable $S$ has the gamma density proportional to $s^{\alpha-1}e^{-\alpha s}$ on $(0,\infty)$, and
\begin{equation}\label{eq:moments}
m_j=\EE S^j=\frac{(\alpha)_j}{\alpha^j}.
\end{equation}
In particular $f$ is a polynomial in $\theta$.

\section{The landscape in the products}\label{sec:products}
Write a polynomial with constant term one as $r_c(s)=1+\sum_{j=1}^{L}c_js^j$ and put $G(c)=\EE[r_c(S)^2]=1+2b^{\top}c+c^{\top}Hc$ with $b_j=m_j$ and $H_{ij}=m_{i+j}$. For a nonzero vector $u\in\RR^L$, $u^{\top}Hu=\EE[(\sum_ju_jS^j)^2]>0$ because $S$ has a positive density, so $H$ is positive definite, $G$ has the unique minimizer $c_*=-H^{-1}b$, and
\begin{equation}\label{eq:quadratic}
G(c)-G(c_*)=(c-c_*)^{\top}H(c-c_*).
\end{equation}
The minimizer is characterized by the normal equations $\EE[S^jr_{c_*}(S)]=0$ for $1\le j\le L$.

\begin{lemma}\label{lem:laguerre}
The minimizing polynomial is $r_*$ of~\eqref{eq:rstar}. It has $L$ simple zeros, all positive, and $G(c_*)=f_*$ as in~\eqref{eq:risk} (see~\cite{PS} for $n=2$ and~\cite{CGPSP} for Gamma measures).
\end{lemma}

\begin{proof}
By~\eqref{eq:moments}, for $1\le j\le L$,
\[
\EE[S^jr_*(S)]=\frac{(\alpha)_j}{\alpha^j}\sum_{h=0}^{L}(-1)^h\binom Lh\frac{(\alpha+j)_h}{(\alpha+1)_h},
\qquad \frac{(\alpha+j)_h}{(\alpha+1)_h}=\frac{(\alpha+1+h)_{j-1}}{(\alpha+1)_{j-1}}.
\]
The right side is a polynomial in $h$ of degree $j-1<L$, so its $L$-th finite difference, which is the alternating sum up to sign, vanishes. Hence $r_*$ satisfies the normal equations and is the minimizer. Equivalently, $r_*$ is orthogonal to all polynomials of degree below $L$ for the weight $s^{\alpha}e^{-\alpha s}$, which identifies it with the rescaled Laguerre polynomial in~\eqref{eq:rstar}; zeros of orthogonal polynomials on $(0,\infty)$ are real, simple and positive~\cite[Theorem~3.3.1]{Szego}.

For the risk, $r_*-1$ lies in the span of $s,\dots,s^L$, to which $r_*$ is orthogonal in $L^2(S)$, so $\EE r_*^2=\EE r_*$. Using $(\alpha)_j/(\alpha+1)_j=\alpha/(\alpha+j)$,
\[
f_*=\sum_{j=0}^{L}(-1)^j\binom Lj\frac{\alpha}{\alpha+j}=\alpha\int_0^1t^{\alpha-1}(1-t)^L\,dt=\alpha B(\alpha,L+1)=\frac{L!}{(\alpha+1)_L}.\qedhere
\]
\end{proof}

The coefficients of $r_a$ are $c_j(a)=(-1)^j\sigma_j(a)$, with $\sigma_j$ the elementary symmetric polynomials, so $F(a)=G(c(a))$. Write $\delta_1,\dots,\delta_L$ for the coordinate vectors of $\RR^L$.

\begin{lemma}\label{lem:Fland}
\begin{enumerate}[label=\textup{(\alph*)}]
\item Every sublevel set of $F$ is compact.
\item Every critical point of $F$ is a global minimizer or has two equal coordinates $a_i=a_j$; in the second case $(\delta_i-\delta_j)^{\top}\Hess F(a)(\delta_i-\delta_j)<0$.
\item The global minimizers of $F$ are the permutations of $(1/\rho_1,\dots,1/\rho_L)$, where $\rho_1,\dots,\rho_L$ are the zeros of $r_*$.
\end{enumerate}
\end{lemma}

\begin{proof}
(a) By~\eqref{eq:quadratic} a sublevel of $F$ bounds every $c_j(a)$, say by $C$. The $a_i$ are the zeros of $z^L+c_1z^{L-1}+\cdots+c_L$, and for $|z|>1+C$ we have $|\sum_jc_jz^{L-j}|<C|z|^L/(|z|-1)<|z|^L$, so all $|a_i|\le1+C$.

(b) If the $a_i$ are distinct, the Jacobian of $a\mapsto c(a)$ is invertible: if it annihilates $u$, then $\sum_iu_i\prod_{h\ne i}(z-a_h)=0$ identically, and evaluating at $z=a_i$ gives $u_i=0$. So at a critical point $a$ of $F$ with distinct coordinates, $c(a)$ is a critical point of $G$; hence $c(a)=c_*$ and $a$ is a global minimizer of $F$ by (c). If $a_i=a_j=\gamma$, put $R(s)=\prod_{l\ne i,j}(1-a_ls)$. Since $(1-(\gamma+t)s)(1-(\gamma-t)s)=(1-\gamma s)^2-t^2s^2$,
\[
F(a+t(\delta_i-\delta_j))=F(a)-2t^2\,\EE\bigl[S^2(1-\gamma S)^2R(S)^2\bigr]+t^4\,\EE\bigl[S^4R(S)^2\bigr],
\]
and both expectations are positive, which gives a second derivative $-4\EE[S^2(1-\gamma S)^2R(S)^2]<0$ at $t=0$.

(c) By Lemma~\ref{lem:laguerre}, $r_*(s)=\prod_i(1-s/\rho_i)$ with distinct $\rho_i>0$. Since $F(a)=G(c(a))\ge G(c_*)$, with equality exactly when $c(a)=c_*$, that is when $r_a=r_*$, the global minimizers of $F$ are the $a$ whose coordinates are, in some order, the $1/\rho_i$.
\end{proof}

\section{The flow on the factors}\label{sec:flow}
Put $g_i=\partial_{a_i}F(a)$ and $h_i=k_i^2+q_i^2$. The flow reads
\begin{equation}\label{eq:flow}
k_i'=-q_ig_i,\qquad q_i'=-k_ig_i,\qquad a_i'=-h_ig_i,\qquad f'=-\|\nabla f\|^2.
\end{equation}
The imbalance $D_i=k_i^2-q_i^2$ is conserved, since $D_i'=-2k_iq_ig_i+2q_ik_ig_i=0$, and
\begin{equation}\label{eq:imbalance}
h_i^2=D_i^2+4a_i^2.
\end{equation}

\begin{lemma}\label{lem:bounded}
Every solution of~\eqref{eq:flow} exists for all $t\ge0$ and stays in a compact set.
\end{lemma}

\begin{proof}
Since $f$ decreases, $a(t)$ stays in a sublevel of $F$, so by Lemma~\ref{lem:Fland}(a) there is a constant $A$, depending on $\theta(0)$, with $|a_i(t)|\le A$. By~\eqref{eq:imbalance}, $k_i^2+q_i^2\le\sqrt{D_i(0)^2+4A^2}$ for all $t$ in the existence interval. The vector field is polynomial, so a solution that stays in a compact set exists for all time.
\end{proof}

Sublevel sets of $f$ itself are not bounded, since $k_i$ can grow while $q_i$ shrinks with $k_iq_i$ fixed; the conserved imbalance is what gives boundedness.

\begin{lemma}\label{lem:saddle}
Every critical point of $f$ that is not a global minimizer is a strict saddle.
\end{lemma}

\begin{proof}
Suppose first that both factors of some layer vanish, $k_i=q_i=0$. Then $a_i=0$ and, with $R_i(s)=\prod_{j\ne i}(1-a_js)$,
\[
g_i=-2\,\EE\bigl[S\,R_i(S)^2\bigr]<0.
\]
Moving only the factors of this layer to $(k_i,q_i)=(t,t)$ changes $a_i$ to $t^2$ and gives the value $f(\theta)-2t^2\EE[SR_i^2]+t^4\EE[S^2R_i^2]$, whose second derivative at $t=0$ is negative. This holds at every point where both factors of some layer vanish.

Suppose now that $(k_i,q_i)\ne(0,0)$ in every layer. Let $J$ denote the Jacobian matrix of the map $\pi(\theta)=(k_iq_i)_i$. It has $i$-th row $(q_i,k_i)$ in the entries of layer $i$ and zero elsewhere, so it has full row rank, and $\nabla f=J^{\top}\nabla F$ forces $\nabla F(a)=0$. If the $a_i$ are distinct, $\theta$ is a global minimizer by Lemma~\ref{lem:Fland}. Otherwise take $a_i=a_j$ and $\xi=\delta_i-\delta_j$, and choose $u$ layer by layer with $\pi(\theta+tu)=a+t\xi$ exactly, for instance by moving $k_l$ by $t\xi_l/q_l$ when $q_l\ne0$ and $q_l$ by $t\xi_l/k_l$ otherwise. Then $f(\theta+tu)=F(a+t\xi)$ and Lemma~\ref{lem:Fland}(b) gives a negative second derivative.
\end{proof}

\begin{lemma}\label{lem:single}
Every solution of~\eqref{eq:flow} converges to a critical point of $f$.
\end{lemma}

\begin{proof}
This is the theorem of \L{}ojasiewicz~\cite{ColdingMinicozzi}; the argument is short. By Lemma~\ref{lem:bounded} and~\eqref{eq:flow}, $f(\theta(t))$ decreases to a limit $f_\infty$ and $\int_0^\infty\|\nabla f(\theta(t))\|^2dt<\infty$. The derivative of $\|\nabla f(\theta(t))\|^2$ is $-2\nabla f^{\top}\Hess f\,\nabla f$, which is bounded on the compact closure of the trajectory, so $\|\nabla f(\theta(t))\|\to0$ and every accumulation point $p$ is critical with $f(p)=f_\infty$. Since $f$ is a polynomial, it is real analytic, and the \L{}ojasiewicz inequality~\cite[(0.2)]{ColdingMinicozzi} gives a ball $U$ around $p$, $C>0$ and $\beta\in(0,1)$ with $|f(\theta)-f_\infty|^{\beta}\le C\|\nabla f(\theta)\|$ on $U$. If $f(\theta(t_0))=f_\infty$ at some time, the trajectory is stationary from then on. Otherwise $\phi(t)=f(\theta(t))-f_\infty>0$ and on $U$
\[
-\frac{d}{dt}\phi^{1-\beta}=(1-\beta)\phi^{-\beta}\|\nabla f\|^2\ge\frac{1-\beta}{C}\|\nabla f\|=\frac{1-\beta}{C}\|\theta'\|.
\]
Let $r$ be the radius of $U$ and choose $T$ with $\|\theta(T)-p\|<r/4$ and $C\phi(T)^{1-\beta}/(1-\beta)<r/4$. Up to a first exit from $U$ after $T$ the trajectory has length below $r/4$, so it never exits, has finite length, and converges; the limit is $p$.
\end{proof}

\begin{proof}[Proof of Theorem~\ref{thm:main}]
Parts (i) and (ii) are Lemmas~\ref{lem:bounded}, \ref{lem:single} and~\ref{lem:saddle}, and (iv) follows from Section~\ref{sec:reduction} and Lemmas~\ref{lem:laguerre} and~\ref{lem:Fland}. For (iii), let $\Sigma$ be the set of strict saddles of $f$. These points also satisfy the definition of a strict saddle in Appendix~A of~\cite{OSS}, which asks the linearized field to have an eigenvalue of each sign. If both factors of layer $i$ vanish, the principal Hessian block of that layer is $\bigl(\begin{smallmatrix}0&g_i\\g_i&0\end{smallmatrix}\bigr)$ with $g_i<0$, so the Hessian has positive as well as negative curvature. Otherwise, for a layer with $q_i\ne0$ we get $\partial_{k_i}^2f=q_i^2\,\partial_{a_i}^2F=2q_i^2\,\EE\bigl[S^2R_i(S)^2\bigr]>0$, with $k_i$ and $q_i$ exchanged if $q_i=0$. Thus $-\Hess f$ has an eigenvalue of each sign at every point of $\Sigma$. The state space is the open set $\RR^{2L}$, the field $-\nabla f$ is $C^1$, and Lemma~\ref{lem:bounded} gives forward completeness. For a gradient flow of a $C^2$ function whose trajectories are defined for all positive time, the set of initial points whose trajectories converge to a point of $\Sigma$ has Lebesgue measure zero~\cite[Appendix~A, Theorem~3 and Corollary~3]{OSS}; there $\Sigma$ may be uncountable. By (i) every trajectory converges to a critical point, and by (ii) a limit outside $\Sigma$ is a global minimizer, so $\mathcal N$ is contained in that null set.
\end{proof}

If $\theta(t)$ converges to a global minimizer, the limiting factor magnitudes are determined by the limiting products and the conserved imbalances: by~\eqref{eq:imbalance},
\[
k_i^2=\tfrac12\Bigl(\sqrt{D_i^2+4a_i^2}+D_i\Bigr),\qquad q_i^2=\tfrac12\Bigl(\sqrt{D_i^2+4a_i^2}-D_i\Bigr),
\]
with $k_i$ and $q_i$ of the same sign, the sign of $k_i(0)+q_i(0)$, since $k_i+q_i$ solves $(k_i+q_i)'=-g_i(k_i+q_i)$ and does not change sign. The products satisfy $a_i'=-h_i\,\partial_{a_i}F(a)$ with $h_i=\sqrt{D_i^2+4a_i^2}$. This is the gradient flow of $F$ for the metric $\operatorname{diag}(h_1^{-1},\dots,h_L^{-1})$, the scalar case of the Riemannian flow of~\cite{BRTW}, and the metric depends on the imbalances: two initial points with the same products and different imbalances give different trajectories of $a$. For $n=2$, $L=2$, $(k_1,q_1)=(1,1)$ and $(k_2,q_2)=(2,1/2)$ both products equal $1$, $h_1=2$, $h_2=17/4$ and $(a_1-a_2)'(0)=99/2$: two layers with equal products and different imbalances separate.

\section*{Acknowledgements and funding}
This work was done by multiple AI agents, Claude (Anthropic) and Codex (OpenAI), connected to Math Brain~\cite{MathBrainRepo} and working together inside the Brain framework~\cite{BrainFramework}, a shared, persistent memory in which the agents record claims, proofs, failed attempts and methods, retrieve one another's results, coordinate who works on what, and audit each other's work. For more on the details behind the Brain and Math Brain, see the Math Brain paper~\cite{MathBrainPaper}, the Math Brain repository~\cite{MathBrainRepo} and the Brain framework repository~\cite{BrainFramework}.

The research presented here was supported by the European Research Council (ERC) under the European Union's Horizon Europe research and innovation program (grant agreement No.~101041711) and by the Simons Foundation as part of the Collaboration on the Mathematical and Scientific Foundations of Deep Learning. It was also supported by Heights Labs, by Convex Nexus Capital, by the Israel Science Foundation (grant number 2258/19), and by the Israel Science Foundation (ISF Grant 4101/25). Computational resources were provided in part by Google Cloud through its education credits program (credit 530550103).

\section*{Declaration of competing interest}
The author holds equity in Heights Labs and in Convex Nexus Capital and has a management role at Convex Nexus Capital. These entities are included in the funding declaration above.

\section*{Declaration of generative AI and AI-assisted technologies}
The paper comes out of Math Brain, an ongoing project on the use of AI systems for research mathematics. The work was done by AI agents, Claude (Anthropic) and Codex (OpenAI), connected to Math Brain~\cite{MathBrainRepo} and working inside the Brain framework~\cite{BrainFramework}, a shared, persistent memory in which the agents record claims, proofs, failed attempts and methods and audit each other's work; the Math Brain paper~\cite{MathBrainPaper} describes both. In November 2026 Codex (OpenAI, GPT-6.1) developed the proofs presented here and wrote the verification script, and Claude Opus~5.5 (Anthropic), working in Claude Code, checked the proofs and wrote the paper. The gradient descent view of attention layers is due to von Oswald et al.~\cite{vonOswald} and Aky\"urek et al.~\cite{Akyurek}. The multi-layer landscape was described by Ahn, Cheng, Daneshmand and Sra~\cite{Ahn}. Convergence of gradient flow was proved for one layer by Zhang, Frei and Bartlett~\cite{ZFB} and for looped models by Gatmiry, Saunshi, Reddi, Jegelka and Kumar~\cite{Gatmiry}, whose value map and mask are used here. Untied attention dynamics were studied by Bordelon, Letey and Pehlevan~\cite{BLP}. Laguerre polynomials as the optimal residual polynomials of gradient methods appear in the work of Pedregosa and Scieur~\cite{PS} and of Cunha et al.~\cite{CGPSP}, and their zeros are treated by classical theory~\cite{Szego}. The conserved imbalance is the balancedness of Arora, Cohen and Hazan~\cite{ACH} and Du, Hu and Lee~\cite{DHL}, and the flow of the products is the scalar case of the Riemannian gradient flow of Bah, Rauhut, Terstiege and Westdickenberg~\cite{BRTW}. The convergence argument uses the \L{}ojasiewicz inequality in the form of Colding and Minicozzi~\cite{ColdingMinicozzi} and the strict-saddle theorem of de~Oliveira, Siami and Sontag~\cite{OSS}. The contribution is the boundedness and convergence of every trajectory, the classification of the critical points that are not global minimizers as strict saddles, and the resulting almost-everywhere convergence for untied factored layers in Theorem~\ref{thm:main}, with the exceptional initializations of Remark~\ref{rem:exceptional}. The author directed the work and is responsible for the content of the paper and for decisions concerning its dissemination.

\section*{Data availability}
The proofs do not depend on computation. The \LaTeX{} source of the paper and a Python script are available at \url{https://github.com/yspennstate/factored-scalar-attention-gradient-flow/tree/v1.0.1}. The script evaluates~\eqref{eq:risk}, \eqref{eq:moments}, \eqref{eq:flow}, \eqref{eq:imbalance}, the Hessians of Lemma~\ref{lem:saddle} and the examples of Remark~\ref{rem:exceptional} and Section~\ref{sec:flow} in exact rational arithmetic for small $n$ and $L$.

\begin{thebibliography}{99}
\bibitem{Ahn} K.~Ahn, X.~Cheng, H.~Daneshmand and S.~Sra, \emph{Transformers learn to implement preconditioned gradient descent for in-context learning}, Advances in Neural Information Processing Systems 36 (2023), arXiv:2306.00297.
\bibitem{ACH} S.~Arora, N.~Cohen and E.~Hazan, \emph{On the optimization of deep networks: implicit acceleration by overparameterization}, Proceedings of Machine Learning Research 80 (2018), arXiv:1802.06509.
\bibitem{Akyurek} E.~Aky\"urek, D.~Schuurmans, J.~Andreas, T.~Ma and D.~Zhou, \emph{What learning algorithm is in-context learning? Investigations with linear models}, International Conference on Learning Representations (2023), arXiv:2211.15661.
\bibitem{BRTW} B.~Bah, H.~Rauhut, U.~Terstiege and M.~Westdickenberg, \emph{Learning deep linear neural networks: Riemannian gradient flows and convergence to global minimizers}, arXiv:1910.05505.
\bibitem{BLP} B.~Bordelon, M.~I.~Letey and C.~Pehlevan, \emph{Theory of scaling laws for in-context regression: depth, width, context and time}, International Conference on Learning Representations (2026), arXiv:2510.01098.
\bibitem{ColdingMinicozzi} T.~H.~Colding and W.~P.~Minicozzi~II, \emph{\L{}ojasiewicz inequalities and applications}, in \emph{Surveys in Differential Geometry 2014: Regularity and Evolution of Nonlinear Equations}, Surveys in Differential Geometry 19, International Press, Somerville, MA, 2015, 63--82, arXiv:1402.5087.
\bibitem{CGPSP} L.~Cunha, G.~Gidel, F.~Pedregosa, D.~Scieur and C.~Paquette, \emph{Only tails matter: average-case universality and robustness in the convex regime}, Proceedings of Machine Learning Research 162 (2022), 4474--4491, arXiv:2206.09901.
\bibitem{DHL} S.~S.~Du, W.~Hu and J.~D.~Lee, \emph{Algorithmic regularization in learning deep homogeneous models: layers are automatically balanced}, Advances in Neural Information Processing Systems 31 (2018), arXiv:1806.00900.
\bibitem{Garg} S.~Garg, D.~Tsipras, P.~Liang and G.~Valiant, \emph{What can transformers learn in-context? A case study of simple function classes}, Advances in Neural Information Processing Systems 35 (2022), arXiv:2208.01066.
\bibitem{Gatmiry} K.~Gatmiry, N.~Saunshi, S.~J.~Reddi, S.~Jegelka and S.~Kumar, \emph{Can looped transformers learn to implement multi-step gradient descent for in-context learning?}, Proceedings of Machine Learning Research 235 (2024), 15130--15152, arXiv:2410.08292.
\bibitem{OSS} A.~C.~B.~de~Oliveira, M.~Siami and E.~D.~Sontag, \emph{Convergence analysis of gradient flow for overparameterized LQR formulations}, Automatica 182 (2025), 112504, arXiv:2408.15456.
\bibitem{PS} F.~Pedregosa and D.~Scieur, \emph{Acceleration through spectral density estimation}, Proceedings of Machine Learning Research 119 (2020), 7553--7562, arXiv:2002.04756.
\bibitem{Szego} G.~Szeg\H{o}, \emph{Orthogonal polynomials}, 4th ed., American Mathematical Society Colloquium Publications 23, Providence, RI, 1975.
\bibitem{vonOswald} J.~von~Oswald, E.~Niklasson, E.~Randazzo, J.~Sacramento, A.~Mordvintsev, A.~Zhmoginov and M.~Vladymyrov, \emph{Transformers learn in-context by gradient descent}, Proceedings of Machine Learning Research 202 (2023), 35151--35174, arXiv:2212.07677.
\bibitem{ZFB} R.~Zhang, S.~Frei and P.~L.~Bartlett, \emph{Trained transformers learn linear models in-context}, Journal of Machine Learning Research 25 (2024), no.~49, arXiv:2306.09927.
\bibitem{MathBrainRepo} Y.~Shmalo, \emph{Math Brain: public snapshot, code and tests}, 2026, \url{https://github.com/yspennstate/math-brain-public/tree/8c9a6bd478c4ec4082347151b7ae16db50c917dd}.
\bibitem{MathBrainPaper} Y.~Shmalo, \emph{Math Brain: shared, auditable memory for multi-agent mathematical research}, October 2026, \url{https://github.com/yspennstate/math-brain-public/tree/main/papers/math-brain-2026-10-24}, doi:10.5281/zenodo.24150837.
\bibitem{BrainFramework} Y.~Shmalo, \emph{Brain Framework}, public software repository, 2026, \url{https://github.com/yspennstate/brain-framework-public/tree/5629e644cdd510173ae97e9aa92642dabd22dc57}.
\end{thebibliography}
\end{document}
